\documentclass[12pt]{article}

\usepackage[a4paper,margin=1in]{geometry}
\usepackage{amsmath,amssymb}
\usepackage{enumitem}
\usepackage{lmodern}
\usepackage[T1]{fontenc}
\usepackage{microtype}

\title{\textbf{MATH6310 Metric Spaces and Topology Coursework 1}}
\author{Sai Vanam}
\date{}

\begin{document}

\maketitle

\section*{Exercise 0.1: Discrete Metric}

Let \(S\) be a set. Define the discrete metric \(d:S\times S\to \mathbb{R}\) by
\[
d(s,t)=
\begin{cases}
0, & s=t,\\
1, & s\neq t.
\end{cases}
\]

We verify that \(d\) satisfies the three metric axioms.

\subsection*{MS1: Identity of Indiscernibles}

By definition,
\[
d(s,t)=0 \iff s=t.
\]

Hence the first metric axiom is satisfied.

\subsection*{MS2: Symmetry}

We show that
\[
d(s,t)=d(t,s).
\]

If \(s=t\), then
\[
d(s,t)=0=d(t,s).
\]

If \(s\neq t\), then \(t\neq s\), and therefore
\[
d(s,t)=1=d(t,s).
\]

Thus,
\[
\boxed{d(s,t)=d(t,s)}
\]
for all \(s,t\in S\), so symmetry holds.

\subsection*{MS3: Triangle Inequality}

We must prove that
\[
d(s,t)\leq d(s,u)+d(u,t)
\]
for all \(s,t,u\in S\).

There are two cases.

\begin{enumerate}[label=\textbf{Case \arabic*:}]
    \item Suppose \(s=t\). Then
    \[
    d(s,t)=0.
    \]
    Since distances are non-negative,
    \[
    0\leq d(s,u)+d(u,t).
    \]
    Hence the triangle inequality holds.

    \item Suppose \(s\neq t\). Then
    \[
    d(s,t)=1.
    \]
    It is impossible for both \(s=u\) and \(u=t\) to hold, since that would imply \(s=t\). Therefore at least one of
    \[
    d(s,u), \qquad d(u,t)
    \]
    must equal \(1\). Hence
    \[
    d(s,u)+d(u,t)\geq 1=d(s,t).
    \]
\end{enumerate}

Therefore,
\[
\boxed{d(s,t)\leq d(s,u)+d(u,t)}.
\]

Hence the discrete metric satisfies all three metric axioms and is therefore a metric on \(S\).

\newpage

\section*{Exercise 0.4: Counting Metric}

Let the counting metric be defined by
\[
d(A,B)=|A\triangle B|,
\]
where \(A\triangle B\) denotes the symmetric difference
\[
A\triangle B=(A\setminus B)\cup(B\setminus A).
\]

We verify the triangle inequality:
\[
d(A,C)\leq d(A,B)+d(B,C).
\]

The key observation is
\[
A\triangle C\subseteq (A\triangle B)\cup(B\triangle C).
\]

To see why, take any \(x\in A\triangle C\). Then \(x\) belongs to exactly one of \(A\) and \(C\).

If \(x\in A\) and \(x\notin C\), then:
\begin{itemize}
    \item if \(x\notin B\), then \(x\in A\triangle B\);
    \item if \(x\in B\), then \(x\in B\triangle C\).
\end{itemize}

Similarly, if \(x\in C\) and \(x\notin A\), then \(x\) must belong to either \(A\triangle B\) or \(B\triangle C\).

Therefore,
\[
A\triangle C\subseteq (A\triangle B)\cup(B\triangle C).
\]

Taking cardinalities gives
\[
|A\triangle C|
\leq
|(A\triangle B)\cup(B\triangle C)|.
\]

Using
\[
|X\cup Y|\leq |X|+|Y|,
\]
we obtain
\[
|(A\triangle B)\cup(B\triangle C)|
\leq
|A\triangle B|+|B\triangle C|.
\]

Hence,
\[
|A\triangle C|
\leq
|A\triangle B|+|B\triangle C|.
\]

Since
\[
d(A,B)=|A\triangle B|,
\]
it follows that
\[
\boxed{d(A,C)\leq d(A,B)+d(B,C)}.
\]

Thus the counting metric satisfies the triangle inequality.

\end{document}
